\section{Parameter-Efficient Fine-Tuning in Continual Learning}
\label{sec:peft}

Parameter-Efficient Fine-Tuning (PEFT), particularly Low-Rank Adaptation (LoRA) \parencite{hu2021loralowrankadaptationlarge}, has become a standard approach to adapt Large Language Models to specialized domains. LoRA operates by freezing the pre-trained weights $W_0 \in \mathbb{R}^{d \times k}$ and representing weight updates $\Delta W$ through the product of two low-rank matrices $B \in \mathbb{R}^{d \times r}$ and $A \in \mathbb{R}^{r \times k}$, where rank $r \ll \min(d, k)$. This formulation reduces the number of trainable parameters up to $10,000$ times, significantly lowering VRAM requirements.

\subsection{Vanilla LoRA and Its Capacity Limits}
Research by Biderman et al. \parencite{biderman2024loralearnsforgets} establishes that LoRA learns less and forgets less compared to full fine-tuning. By constraining updates to a low-rank subspace, LoRA acts as a geometric regularizer that prevents the weights from drifting too far from the pre-trained distribution, thereby preserving generation diversity and general capabilities.

However, Vanilla LoRA faces capacity limits: the intrinsic rank of the parameter update $\Delta W$ in full-parameter continual pre-training (CPT) on a new domain is typically around $2,000$. Consequently, low-rank adapters with conventional ranks (such as $r \le 256$) lack the capacity to absorb massive factual knowledge during domain pre-training. LoRA is therefore best suited for instruction fine-tuning (IFT) to adapt task behaviors rather than memorizing domain-specific facts.

\subsection{Subspace Isolation: O-LoRA and N-LoRA}
Fine-tuning sequential tasks on a single LoRA adapter often leads to parameter overwriting. To address this, O-LoRA \parencite{wang-etal-2023-orthogonal, repec:axf:aidtaa:v:3:y:2026:i:1:p:52-61} projects gradient updates of new tasks onto a subspace orthogonal to the LoRA parameters learned on previous tasks. While O-LoRA minimizes cross-task interference, it remains highly sensitive to hyperparameters like rank $r$ and scaling factor $\alpha$, which can lead to performance drops in complex reasoning tasks.

N-LoRA \parencite{yang2024parametercollisionhinderingcontinual} identifies that orthogonal subspaces do not fully eliminate parameter collisions at the scalar level. It introduces an $\ell_1$ sparsity constraint to force LoRA parameters to reside in sparse, disjoint regions. N-LoRA reduces parameter collisions $58.1$ times compared to O-LoRA, providing stronger protection against catastrophic forgetting.

\subsection{Architectural Extensions: MoE-LoRA and Prompts}
To scale to multiple tasks, LoRA has been combined with Mixture of Experts (MoE). The SLIM framework \parencite{han-etal-2025-slim} applies soft routing between LoRA experts and identity layers, allowing general queries to bypass the adapters entirely. MoLA \parencite{gao2024higherlayersneedlora} allocates LoRA experts using an inverted triangle layout: fewer experts are allocated in lower layers (which extract stable syntactic features), while more experts are assigned in higher layers (which extract semantic features vulnerable to interference). Additionally, I-LoRA \parencite{ren2024analyzingreducingcatastrophicforgetting} connects adjacent task minima through linear weight interpolation, while Progressive Prompts \parencite{razdaibiedina2023progressivepromptscontinuallearning} concatenates task-specific soft prompt tokens while keeping base weights frozen.

Despite these advances, managing multiple adapters or routing mechanisms increases GPU memory overhead and runtime complexity. Furthermore, PEFT methods still modify parameters permanently. These limitations motivate a shift toward semi-parametric architectures, where non-parametric external memory (RAG) is utilized for factual knowledge storage.
